% =============================================================================
% ANNEXE M : TRIANGULATION EXTERNE A2 -- EXTRAIT DE COURRIEL DE CONVERGENCE NDT
% =============================================================================

\chapter{Triangulation externe : extrait de courriel attestant la convergence sur la NDT}
\label{annex:triangulation-a2-mail}

\annexefiche
  {Courriel professionnel produit chez TWIICE en 2026, postérieurement à la réunion d'alignement sur la NDT-SW-001 (horodatage exact caviardé)}
  {Message interne adressé par l'équipe Qualité/Affaires Réglementaires (QA/RA) aux parties prenantes de la réunion (R\&D, Software, l'auteur du portfolio)}
  {Triangulation externe de la \textbf{Preuve~1 de A2} (note de décision technique, annexe~\ref{annex:ndt-sw-001}) : un tiers atteste par écrit que la convergence a été obtenue et reconnaît la NDT comme référence partagée. Mobilisé en section A2.2 du chapitre \textit{Situations professionnalisantes}.}
  {Caviardage : expéditeur, destinataires nominatifs, fonctions précises, horodatage, références internes de projet et de procédure. Le contenu décisionnel est conservé intégralement.}

% : Note de mise en contexte (portfolio) -------------------------------------
\begin{tcolorbox}[colback=blue!5, colframe=blue!40!black, boxrule=0.5pt, arc=3pt,
  left=8pt, right=8pt, top=6pt, bottom=6pt]
  \small\itshape
  \textbf{Note du portfolio.} Cet artefact est un courriel professionnel réel,
  envoyé par l'équipe Qualité/Affaires Réglementaires (QA/RA) de TWIICE à
  l'issue de la réunion d'alignement portant sur la note de décision technique
  NDT-SW-001 (cf.~annexe~\ref{annex:ndt-sw-001}).

  Il atteste, par la parole écrite d'un tiers institutionnel distinct de
  l'auteur du portfolio, que la convergence revendiquée en section A2.2 a
  effectivement été obtenue et a été \emph{reconnue comme telle} par la
  fonction qualité, indépendante hiérarchiquement de l'équipe logicielle.

  Conformément aux règles de confidentialité de l'entreprise, l'expéditeur,
  les destinataires nominatifs, les fonctions précises et les références
  internes de projet ont été \textbf{caviardés}. Le contenu décisionnel
  (alignements, actions, points ouverts, prochaines étapes) est en revanche
  conservé \textbf{intégralement} : c'est lui qui porte la valeur probante.

  Au sens de Tardif, ce courriel transforme la triangulation externe pour A2
  d'un statut \emph{disponible mais non mobilisée} en un statut
  \emph{partiellement insérée et vérifiable}, sans avoir nécessité la
  divulgation d'informations confidentielles.
\end{tcolorbox}

\bigskip

% : Reproduction caviardee du courriel ---------------------------------------
\begin{tcolorbox}[enhanced, breakable, colback=gray!3, colframe=gray!55!black,
  boxrule=0.4pt, arc=2pt,
  left=12pt, right=12pt, top=10pt, bottom=10pt,
  title={\small\bfseries Courriel reproduit (caviardé)},
  fonttitle=\small\bfseries, coltitle=white,
  boxed title style={colback=gray!55!black, arc=2pt, boxrule=0pt}]

\small
\begin{tabular}{@{}p{2.4cm}@{\hspace{0.6em}}p{\dimexpr\linewidth-3cm\relax}@{}}
\textbf{De} & \caviarder[5.5cm] \hfill\textit{(QA/RA, TWIICE)} \\[2pt]
\textbf{À}  & \caviarder[5.5cm] \hfill\textit{(R\&D, Software, auteur)} \\[2pt]
\textbf{Date} & \caviarder[3cm] \\[2pt]
\textbf{Objet} & Compte-rendu de la réunion d'alignement NDT~--- décisions et actions \\
\end{tabular}

\medskip
\hrule
\medskip

Bonjour,

\medskip
Merci pour la réunion. Voici le résumé des points clés.

\medskip
\textbf{Décisions / alignements}
\begin{itemize}
  \item Convergence confirmée sur la NDT comme référence du besoin et point
        d'ancrage de la traçabilité.
  \item Accord pour aligner les exigences système et software sur la NDT,
        sans interprétations implicites. Toute divergence devra être justifiée
        et documentée.
\end{itemize}

\textbf{Actions}
\begin{itemize}
  \item QA/RA met à jour la matrice de traçabilité (NDT~$\rightarrow$~exigences
        $\rightarrow$~vérifications) pour refléter l'alignement.
  \item R\&D et Software relisent les passages encore ambigus et proposent une
        formulation finale.
\end{itemize}

\textbf{Points ouverts}
\begin{itemize}
  \item Arbitrage sur 2~formulations jugées ambiguës.
  \item Définition d'une date de gel (\emph{baseline}) de la NDT pour stabiliser
        la traçabilité.
\end{itemize}

\textbf{Prochaines étapes}
\begin{itemize}
  \item Mini-revue de clôture pour valider les deux formulations et acter la
        \emph{baseline}.
  \item Ensuite, lancement des mises à jour qualité associées.
\end{itemize}

\medskip
Cordialement,

\caviarder[4cm] \\
\caviarder[5.5cm] \hfill\textit{(QA/RA, TWIICE)}

\end{tcolorbox}

\bigskip

% : Lecture probante ---------------------------------------------------------
\section*{Lecture probante au sens de Tardif}

Quatre éléments du courriel constituent la valeur de triangulation~:

\begin{enumerate}
  \item \textbf{Émetteur tiers.} Le message est envoyé par la fonction QA/RA, hiérarchiquement distincte de l'équipe logicielle et institutionnellement responsable de la traçabilité IEC~62304. Sa parole engage l'organisation, pas l'auteur du portfolio.
  \item \textbf{Reconnaissance explicite de la convergence.} La formulation \emph{« Convergence confirmée sur la NDT comme référence du besoin et point d'ancrage de la traçabilité »} atteste, par un tiers, l'effectivité du résultat revendiqué en section A2.2 du portfolio. C'est la traduction écrite, postérieure à la réunion, du même fait que la NDT-SW-001 documente \emph{ex ante}.
  \item \textbf{Conséquences opérationnelles.} La distribution d'actions à QA/RA (matrice de traçabilité), R\&D et Software, ainsi que la planification d'une \emph{baseline}, montrent que la convergence n'est pas restée déclarative~: elle a déclenché des engagements de production documentaire chez d'autres équipes.
  \item \textbf{Lucidité sur les points ouverts.} La mention de deux formulations encore ambiguës à arbitrer témoigne de l'honnêteté du compte-rendu~: ce n'est pas un message de clôture triomphale mais un point d'étape équilibré, ce qui renforce sa valeur probante.
\end{enumerate}

\medskip
\noindent\textbf{Limite assumée.} Ce courriel est l'unique élément de triangulation externe matériellement inséré dans le portfolio. Les éléments équivalents pour B4 (historique GitLab, validations qualité du SDP) et C2 (feuille de route contresignée par la direction technique) restent \emph{disponibles sur demande motivée} mais non insérés (cf.~encart \emph{Triangulation externe} en tête du chapitre \textit{Situations professionnalisantes}). L'insertion d'un seul artefact suffit néanmoins à transformer le statut global de la triangulation externe d'une posture \emph{déclarative} à une posture \emph{partiellement insérée}, conformément à l'exigence Tardif de vérifiabilité par un tiers.

\clearpage

